% GENERATED FILE. Edit canonical lessons, then run npm run generate:books.
% canonical-source-hash: fnv1a64:87a25aba7fa207fd
% canonical-lessons: AR-C241-muhadhdhab, AR-C241-sabur, AR-C241-mutawadi, AR-C241-mudhik, AR-C241-jadd

\chapter{Polite, Patient, Modest, Funny, Serious}
\label{ch:ar-polite-patient-modest-funny-serious}
\hlchaptermodality{241}

\begin{chapteropening}
\textbf{By the end of this chapter:} \emph{I can say polite, patient, modest, funny and serious.}

Complete the last lesson of chapter 241: I can say polite, patient, modest, funny and serious.
\end{chapteropening}

\section[\texorpdfstring{\ar{مهذب} (muhadhdhab)}{مهذب (muhadhdhab)}]{\ar{مهذب} (muhadhdhab) — polite}

\label{lesson:AR-C241-muhadhdhab}



\begin{quote}
Before the new one: say the Arabic for suitable, then the Arabic for excellent.
\end{quote}

\subsection*{You'll want to know: \ar{مهذب}}
\textbf{\ar{مهذب}} — \emph{muhadhdhab} — \textquotedblleft{}polite\textquotedblright{}.

\begin{cousinweb}
\textbf{\ar{مهذب}} (\emph{muhadhdhab}) is built on the native Arabic root \emph{h-dh-b}, \textquotedblleft{}to prune, to refine\textquotedblright{} (Lane). \emph{Muhadhdhab} is the Form II passive participle, \textquotedblleft{}refined\textquotedblright{}, hence \textquotedblleft{}polite\textquotedblright{}.
\end{cousinweb}

\begin{grammarlens}[title={What you've built}]
One more word for this chapter.
\end{grammarlens}

\subsection*{Your turn}
\begin{itemize}
\raggedright
  \item \emph{Say it:} \emph{muhadhdhab}
  \item \emph{Say it:} \emph{muhadhdhab}, once more
  \item \emph{Say it:} say \emph{muhadhdhab}
  \item \emph{Recall:} say the Arabic for suitable, then the Arabic for excellent, then say \emph{muhadhdhab} again
\end{itemize}

\begin{tcolorbox}[breakable,colback=teal!4,colframe=teal!35!black,title={Before you move on}]
What does \ar{مهذب} mean? (\textquotedblleft{}Polite\textquotedblright{}.) Say it once more.
\end{tcolorbox}
\section[\texorpdfstring{\ar{صبور} (ṣabūr)}{صبور (ṣabūr)}]{\ar{صبور} (ṣabūr) — patient}

\label{lesson:AR-C241-sabur}



\begin{quote}
Before the new one: say the Arabic for excellent, then the Arabic for polite.
\end{quote}

\subsection*{You'll want to know: \ar{صبور}}
\textbf{\ar{صبور}} — \emph{ṣabūr} — \textquotedblleft{}patient\textquotedblright{}.

\begin{cousinweb}
\textbf{\ar{صبور}} (\emph{ṣabūr}) is built on the native Arabic root \emph{ṣ-b-r}, \textquotedblleft{}to bind; to be patient\textquotedblright{} (Lane). \emph{Ṣabūr} is the \emph{faʿūl} adjective, \textquotedblleft{}patient\textquotedblright{}.
\end{cousinweb}

\begin{grammarlens}[title={What you've built}]
One more word for this chapter.
\end{grammarlens}

\subsection*{Your turn}
\begin{itemize}
\raggedright
  \item \emph{Say it:} \emph{ṣabūr}
  \item \emph{Say it:} \emph{ṣabūr}, once more
  \item \emph{Say it:} say \emph{ṣabūr}
  \item \emph{Recall:} say the Arabic for excellent, then the Arabic for polite, then say \emph{ṣabūr} again
\end{itemize}

\begin{tcolorbox}[breakable,colback=teal!4,colframe=teal!35!black,title={Before you move on}]
What does \ar{صبور} mean? (\textquotedblleft{}Patient\textquotedblright{}.) Say it once more.
\end{tcolorbox}
\section[\texorpdfstring{\ar{متواضع} (mutawāḍiʿ)}{متواضع (mutawāḍiʿ)}]{\ar{متواضع} (mutawāḍiʿ) — modest, humble}

\label{lesson:AR-C241-mutawadi}



\begin{quote}
Before the new one: say the Arabic for polite, then the Arabic for patient.
\end{quote}

\subsection*{You'll want to know: \ar{متواضع}}
\textbf{\ar{متواضع}} — \emph{mutawāḍiʿ} — \textquotedblleft{}modest, humble\textquotedblright{}.

\begin{cousinweb}
\textbf{\ar{متواضع}} (\emph{mutawāḍiʿ}) is built on the native Arabic root \emph{w-ḍ-ʿ}, \textquotedblleft{}to put down, to lay\textquotedblright{} (Lane). \emph{Mutawāḍiʿ} is the Form VI participle, \textquotedblleft{}lowering oneself\textquotedblright{}, hence \textquotedblleft{}humble\textquotedblright{}.
\end{cousinweb}

\begin{grammarlens}[title={What you've built}]
One more word for this chapter.
\end{grammarlens}

\subsection*{Your turn}
\begin{itemize}
\raggedright
  \item \emph{Say it:} \emph{mutawāḍiʿ}
  \item \emph{Say it:} \emph{mutawāḍiʿ}, once more
  \item \emph{Say it:} say \emph{mutawāḍiʿ}
  \item \emph{Recall:} say the Arabic for polite, then the Arabic for patient, then say \emph{mutawāḍiʿ} again
\end{itemize}

\begin{tcolorbox}[breakable,colback=teal!4,colframe=teal!35!black,title={Before you move on}]
What does \ar{متواضع} mean? (\textquotedblleft{}Modest, humble\textquotedblright{}.) Say it once more.
\end{tcolorbox}
\section[\texorpdfstring{\ar{مضحك} (muḍḥik)}{مضحك (muḍḥik)}]{\ar{مضحك} (muḍḥik) — funny}

\label{lesson:AR-C241-mudhik}



\begin{quote}
Before the new one: say the Arabic for patient, then the Arabic for modest.
\end{quote}

\subsection*{You'll want to know: \ar{مضحك}}
\textbf{\ar{مضحك}} — \emph{muḍḥik} — \textquotedblleft{}funny\textquotedblright{}.

\begin{cousinweb}
\textbf{\ar{مضحك}} (\emph{muḍḥik}) is built on the native Arabic root \emph{ḍ-ḥ-k}, \textquotedblleft{}to laugh\textquotedblright{} (Lane). \emph{Muḍḥik} is the Form IV participle, \textquotedblleft{}making laugh\textquotedblright{}.
\end{cousinweb}

\begin{grammarlens}[title={What you've built}]
One more word for this chapter.
\end{grammarlens}

\subsection*{Your turn}
\begin{itemize}
\raggedright
  \item \emph{Say it:} \emph{muḍḥik}
  \item \emph{Say it:} \emph{muḍḥik}, once more
  \item \emph{Say it:} say \emph{muḍḥik}
  \item \emph{Recall:} say the Arabic for patient, then the Arabic for modest, then say \emph{muḍḥik} again
\end{itemize}

\begin{tcolorbox}[breakable,colback=teal!4,colframe=teal!35!black,title={Before you move on}]
What does \ar{مضحك} mean? (\textquotedblleft{}Funny\textquotedblright{}.) Say it once more.
\end{tcolorbox}
\section[\texorpdfstring{\ar{جاد} (jādd)}{جاد (jādd)}]{\ar{جاد} (jādd) — serious}

\label{lesson:AR-C241-jadd}



\begin{quote}
Before the new one: say the Arabic for modest, then the Arabic for funny.
\end{quote}

\subsection*{You'll want to know: \ar{جاد}}
\textbf{\ar{جاد}} — \emph{jādd} — \textquotedblleft{}serious\textquotedblright{}.

\begin{cousinweb}
\textbf{\ar{جاد}} (\emph{jādd}) is built on the native Arabic root \emph{j-d-d}, \textquotedblleft{}to cut; to be in earnest\textquotedblright{} (Lane). \emph{Jādd} is the active participle, \textquotedblleft{}in earnest\textquotedblright{}, hence \textquotedblleft{}serious\textquotedblright{}.
\end{cousinweb}

\begin{grammarlens}[title={What you've built}]
One more word for this chapter.
\end{grammarlens}

\subsection*{Your turn}
\begin{itemize}
\raggedright
  \item \emph{Say it:} \emph{jādd}
  \item \emph{Say it:} \emph{jādd}, once more
  \item \emph{Say it:} say \emph{jādd}
  \item \emph{Recall:} say the Arabic for modest, then the Arabic for funny, then say \emph{jādd} again
\end{itemize}

\begin{tcolorbox}[breakable,colback=teal!4,colframe=teal!35!black,title={Before you move on}]
What does \ar{جاد} mean? (\textquotedblleft{}Serious\textquotedblright{}.) Say it once more.
\end{tcolorbox}
